\chapter{Introduction}
\label{chapter:introduction}

\section{Motivation}

Language has become a common interface between people and multimodal machine learning systems. Modern models can process information across text, images, audio, and video~\citep{han2024onellm,radford2021learning,liu2023visual}, yet many systems use words to connect these modalities. In text-to-image generation, users describe the desired result with a prompt~\citep{rombach2022high}; in vision-language models, users ask questions about images and receive answers in language~\citep{liu2023visual}; and in voice assistants, spoken instructions are translated into actions. This design is natural and useful, but it raises a basic question: \emph{Is this the most effective and efficient way?}

Language is powerful, but it can also be a narrow channel. Expression unfolds sequentially, so a person may have to compress a rich intention into a short prompt and then refine it through several rounds of feedback. Some experiences are even harder to verbalize. The abstract structure and affective intensity of Wassily Kandinsky's \emph{Composition VII}, for example, can be perceived as a whole but are difficult to capture precisely in words. Together, these limitations motivate a complementary channel through which some percepts, intentions, or internal states can be communicated without deliberate verbalization.

The same limitation can appear in the opposite direction, when a machine communicates with a person. In patient education, an answer is useful only if a lay reader can understand it without losing essential information. The models evaluated in \repqa, introduced in Chapter~\ref{chapter:repqa}, often miss a requested reading level, while interventions that improve readability can omit part of the original content.

Motivated by these observations, this dissertation develops machine learning methods for a more direct interface between mind and world. These methods model neural signals, connect them to semantic and perceptual representations, and investigate how action meaning can be learned from neural activity and how embodied systems can realize specified goals. Language may still serve as an internal representation or output format; the goal is to reduce the need for a person to formulate every intention as a verbal instruction.

Such an interface could matter in three ways. First, models that identify meaningful patterns in neural activity can help reveal how distributed brain systems support perception, cognition, and behavior. They may also support clinical analysis by characterizing disease-related network changes, predicting treatment response, and evaluating rehabilitation-related neural changes, as investigated respectively in Alzheimer's disease, depression, and multiple sclerosis~\citep{li2015toward,siegle2006use,cerasa2013computer}. Second, neural decoding could make aspects of internal thought, mental imagery, and dreams available in a form that can be inspected and shared. Text could provide a readable account of decoded semantic content, while images could convey visual appearance and spatial relationships. For example, externalizing an imagined scene could provide a starting point for creative work, and recovering aspects of dream imagery could offer material for studying subjective experience. Third, decoded neural signals could provide commands or goals to external systems. Speech neuroprostheses and a neurally controlled robotic arm already show how neural signals can support communication and device control for people with speech or motor impairments~\citep{card2024accurate,hochberg2012reach}; in the longer term, decoded intentions could also guide other assistive devices and embodied agents.

Realizing this interface raises three fundamental challenges. (1) \emph{Understanding the brain} requires a two-way exchange: machine learning must discover neural structure that is not specified in advance, while principles of neural organization can guide the analysis and design of machine learning models. (2) \emph{Decoding brain states} requires identifying which candidate representation is well aligned with neural activity and learning that mapping across subjects and tasks from limited paired data to recover content as text or images. (3) \emph{Enabling interaction} moves from action meaning to goal realization: how can neural representations acquire meaning through actions and their consequences, and how can a person specify a desired outcome while an embodied agent determines how to achieve it effectively and efficiently?

\section{Contributions}

This dissertation addresses these challenges through studies of neural modeling, semantic and perceptual representations, and embodied interaction. Across these studies, each model exploits structure specific to its problem: brain connectivity and anatomy, objects and relations in decoded scenes, and temporal persistence or graph topology in control and search. The main contributions are grouped as follows:

\begin{itemize}
    \item \textbf{Connecting neural organization and model design.} I use machine learning to identify candidate patterns of brain organization and draw on neuroscience to study the organization of language models. \hybrid{} learns candidate outcome-relevant, high-order relationships among brain regions instead of restricting analysis to pairwise connections. Small-World Allocation (SWA) uses the functional organization of attention heads to guide pruning, translating a principle studied in brain networks into a method for preserving language model performance (Chapters~\ref{chapter:hybrid} and~\ref{chapter:swa}).

    \item \textbf{Building a language bridge from fMRI alignment to disentangled visual generation.} The argument begins with \prism{} I, whose comparison of candidate representation spaces motivates language as an alignment target for fMRI. \mindllm{} then learns fMRI-to-language alignment across subjects and tasks, combining a neuroscience-informed encoder with Brain Instruction Tuning: the encoder incorporates voxel location and brain parcellations, while instruction tuning expands supervision through existing annotations linked by shared stimulus images. \prism{} II examines how language can support visual reconstruction by structuring the representation into objects, attributes, and relations and composing scenes object by object. This object-disentangled space motivates \dedit{}, which strengthens the generator's item-level disentanglement through learned per-item prompts and grouped cross-attention, enabling broader visual control. These contributions connect space selection, alignment learning, language-mediated reconstruction, and a more expressive disentangled generator (Chapters~\ref{chapter:prism}--\ref{chapter:dedit}).

    \item \textbf{Moving from action meaning to goal-directed behavior.} I study how intention can become behavior, progressing from the action level to the goal level. The Mind World Model (\mwm{}) jointly learns a mind model that encodes EEG as latent actions and a world model that predicts their visual consequences in synchronized egocentric video, connecting brain signals to action intention without action annotations. Completing a task often calls for specifying the desired outcome rather than each action, motivating two complementary approaches to goal realization. Their interfaces are a goal image for \phil{} and a language instruction for \dysta{}, corresponding to the images and text decoded in Part~\ref{part:bridge}. \phil{} learns graph-search heuristics for a target state representing the outcome specified by the image. \dysta{} generates actions from the language instruction and multi-frame visual history, using static--dynamic disentanglement to improve task success and inference efficiency. These independently evaluated studies connect the research questions of learning action meaning and acting toward a desired result (Chapters~\ref{chapter:mwm}--\ref{chapter:dysta}).
\end{itemize}

My long-term goal is to develop an additional communication channel between people and machines, especially when speech is impossible, too slow, or unable to preserve the relevant content. The present work recovers only task-dependent aspects of neural activity and does not yet integrate neural decoding with robotic control. It develops the main components separately: neural modeling, representation learning, control, and planning.

\section{Thesis Outline}

The dissertation begins with a prologue and then proceeds in three parts: the foundation for modeling brain signals, the bridge from signals to representations, and the study of action meaning and goal realization for embodied interaction.

\textbf{Prologue.} Chapter~\ref{chapter:repqa} introduces \repqa, a benchmark for evaluating readability in patient education question answering. The chapter asks whether large language models can write at a requested reading level, infer the level required by a question, and improve readability without sacrificing informativeness. The results show that the models often fail to meet requested reading levels and that readability-oriented post-training can discard useful content. This chapter provides a concrete example of information being lost in the machine-to-human direction and motivates the search for complementary interfaces.

\textsc{Part~\ref{part:foundation}} examines the relationship between neural organization and machine learning. Its two chapters discover high-order structure in brain signals and apply a principle from neuroscience to language model analysis and pruning.

Chapter~\ref{chapter:hybrid} introduces \hybrid{}, a method for identifying candidate high-order relationships among brain regions. Existing approaches usually represent functional connectivity with pairwise edges or select higher-order candidates before learning, which can miss distributed patterns that matter for cognition. \hybrid{} instead represents these relationships as weighted hyperedges and learns them end to end through a maximally-informative-minimally-redundant objective and a multi-head drop-bottleneck. The learned hyperedges improve phenotypic prediction, while post-hoc analyses suggest that higher-degree hyperedges carry useful phenotypic information.

Chapter~\ref{chapter:swa} studies whether the small-world organization observed in functional brain networks also provides insight into language models. It constructs functional graphs whose nodes are attention heads and whose edges reflect activation similarity. Across the evaluated models, stronger small-world organization was associated with better fluid reasoning performance. SWA uses scores for within-community and cross-community connectivity to allocate sparsity across heads and layers, preserving this organization and improving performance under pruning.

\textsc{Part~\ref{part:bridge}} develops the bridge from brain signals to visual content through a connected argument: \prism{} I informs the choice of alignment space; \mindllm{} learns fMRI-to-language alignment; \prism{} II uses language for visual reconstruction; and the object-disentangled structure of that representation motivates a more expressive disentangled generator in \dedit{}. The two \prism{} contributions remain together in Chapter~\ref{chapter:prism}, while the exposition below follows their roles in this argument.

\textbf{\prism{} I: choosing the alignment space.} Chapter~\ref{chapter:prism} compares representations from vision, language, and vision-language models as targets for fMRI decoding. Among the evaluated alternatives, the experiments found stronger alignment with language-model semantic representations. This finding motivates the choice of language as the space in which to learn a versatile neural decoder.

\textbf{\mindllm{}: learning fMRI-to-language alignment.} Chapter~\ref{chapter:mindllm} introduces \mindllm{}, a subject-agnostic and versatile model for fMRI-to-text decoding. Its neuroscience-informed encoder aggregates variable voxel sets using learnable queries, keys that encode voxel position and brain parcellations, and values carrying measured fMRI activity. Brain Instruction Tuning links existing recordings to annotations from fifteen task datasets through shared stimulus images, yielding 980{,}610 conversations spanning perception, memory and knowledge, language and symbols, and complex reasoning. Together, these components address variation across subjects and scarcity of task-specific supervision, learning alignment to language without new scanning for each task. Experiments evaluated captioning, broader question answering, adaptation to new tasks, and generalization to an unseen subject.

\textbf{\prism{} II: reconstructing visual content through language.} The reconstruction contribution in Chapter~\ref{chapter:prism} asks how the language space can support visual reconstruction. A single scene description can entangle objects with their attributes. \prism{} structures the representation into objects, attributes, and relations and generates the scene object by object, improving reconstruction quality and semantic faithfulness. This object-disentangled space makes preservation of item-level structure a requirement for the generator.

\textbf{\dedit{}: strengthening the disentangled generator.} Chapter~\ref{chapter:dedit} develops the generative side beyond the object-centric construction used in \prism{} II. \dedit{} assigns a separate learned prompt to each image item and introduces grouped cross-attention so that each region is controlled by its associated prompt. This strengthens item-level disentanglement and supports a broader set of visual controls: text-, reference-image-, and mask-guided editing, as well as item removal, within one framework. Its image-editing evaluations establish this generative capability; connecting the generator to decoded neural representations remains future work.

\textsc{Part~\ref{part:embodiment}} asks how intention can become behavior in the world, moving from action meaning to goal realization. At the action level, \mwm{} grounds neural representations through joint learning of a mind model and a world model that predicts visual consequences. The question then shifts to the goal level: how can a person specify a desired result rather than each action? \phil{} and \dysta{} address this question through complementary interfaces: image-specified goals for explicit graph search and language instructions for direct action generation from visual history. These correspond to the images and text decoded in Part~\ref{part:bridge}. The progression organizes the research questions of three independently evaluated models.

Chapter~\ref{chapter:mwm} introduces the Mind World Model (\mwm{}), which connects brain signals to action intention by learning motor EEG representations through the visual consequences of movement, without action annotations. Supervised motor decoding relies on costly cue-based trials and coarse categorical labels, while existing EEG foundation models pretrain with EEG-only objectives that do not directly capture how movement unfolds. \mwm{} treats the EEG recorded up to the current frame of an egocentric video as the latent action of a visual world model, which uses it to predict the representation of a future frame. Its Lateralization-Aware Transformer separates bilateral from side-specific neural activity and exchanges only the side-specific component across hemispheres when the video is mirrored. In linear-probe evaluations, the pretrained mind model outperformed EEG foundation models on all eleven motor-decoding tasks, and the pretrained world model transferred to future prediction conditioned on independently measured IMU signals.

Chapter~\ref{chapter:phil} shifts from action meaning to goal-directed planning, studying explicit graph search as a capability for realizing an image-specified outcome. In the dissertation's embodiment framing, a goal image depicts the desired outcome, which must be grounded in a goal state and its features. Given a state graph, a start state, and that goal state, \phil{} learns graph-search heuristics that use the structure revealed during search. It maintains a digest of the graph explored so far and predicts node-to-goal distances that guide the next expansion. Experiments on grids, graphs from several domains, and simulated indoor UAV path planning show that the learned heuristic reduces search effort. These experiments evaluate supplied graph-and-goal planning problems.

Chapter~\ref{chapter:dysta} studies a complementary policy-based approach at the goal level. \dysta{} is a vision-language-action model that generates actions directly from a language instruction and multi-frame visual history. Its goal interface corresponds to the text output of Part~\ref{part:bridge}, providing a complementary route to \phil{}'s image-specified goal. Effective behavior often depends on earlier observations, but standard models either observe only the current frame or incur a rapidly growing context and repeated computation when past frames are included. \dysta{} separates visual tokens according to their temporal persistence, retains one copy of static information across time, and uses a learned recache gate to decide when cached tokens should be refreshed. This design improves performance on tasks that depend on earlier observations while reducing both context length and inference latency, making the policy more effective and efficient.
